\documentclass[11pt,a4paper]{article}
\usepackage[utf8]{inputenc}
\usepackage[T1]{fontenc}
\usepackage[spanish]{babel}
\usepackage{amsmath,amssymb}
\usepackage[margin=2.5cm]{geometry}

% Compilar desde esta carpeta: pdflatex propuesta_reactivos_MATD124.tex
% Cambiar a \solucionesfalse para generar la versión para estudiantes.
\newif\ifsoluciones
\solucionestrue
\newenvironment{alternativas}{%
  \begin{enumerate}\renewcommand{\labelenumi}{\Alph{enumi})}%
  \setlength{\itemsep}{3pt}\setlength{\parskip}{0pt}%
}{\end{enumerate}}
\newcommand{\reactivo}[2]{\subsection*{Reactivo #1}\noindent #2\par}
\setlength{\parindent}{0pt}
\setlength{\parskip}{6pt}

\begin{document}
\begin{center}
  {\Large\bfseries Propuesta de reactivos de Cálculo Diferencial}\par
  {\large MATD124 --- Examen de media carrera}
\end{center}


\section*{Instrucciones}
Seleccione una sola respuesta en cada pregunta. Cada reactivo tiene cuatro
alternativas y únicamente una es correcta. Salvo indicación contraria, las
variables y funciones son reales; $\mathbb{N}=\{1,2,3,\ldots\}$.

\reactivo{1}{Sea $S=\{1-1/n:n\in\mathbb{N}\}$. ¿Cuál es la afirmación correcta?}
\begin{alternativas}
  \item $\inf S=0$ y $\max S=1$.
  \item $\min S=0$, $\sup S=1$ y $S$ no tiene máximo.
  \item $\inf S=-1$ y $\sup S=1$.
  \item $\min S=0$ y $\sup S=0$.
\end{alternativas}

\reactivo{2}{Para la sucesión $a_n=(3n^2+1)/(2n^2+n)$, ¿cuál es el valor de
$\displaystyle\lim_{n\to\infty}a_n$?}
\begin{alternativas}
  \item $0$.
  \item $1$.
  \item $3$.
  \item $\dfrac{3}{2}$.
\end{alternativas}

\reactivo{3}{¿Cuál de las siguientes condiciones define que una sucesión
real $(a_n)$ es de Cauchy?}
\begin{alternativas}
  \item Para todo $\varepsilon>0$, existe $N\in\mathbb{N}$ tal que
  $|a_n-a_m|<\varepsilon$ para todos $n,m\geq N$.
  \item Existe $\varepsilon>0$ tal que $|a_n-a_m|<\varepsilon$ para todos
  $n,m\in\mathbb{N}$.
  \item Para todo $\varepsilon>0$, existe $N\in\mathbb{N}$ tal que
  $|a_n|<\varepsilon$ para todo $n\geq N$.
  \item Existe $N\in\mathbb{N}$ tal que $a_n=a_m$ para todos $n,m\geq N$.
\end{alternativas}

\reactivo{4}{¿Cuál es el valor de
$\displaystyle\lim_{x\to 0}\frac{e^{2x}-1}{x}$?}
\begin{alternativas}
  \item $0$.
  \item $1$.
  \item $2$.
  \item El límite no existe.
\end{alternativas}

\reactivo{5}{Sea
\[
f(x)=\begin{cases}
\dfrac{x^2-4}{x-2},&x\neq 2,\\
k,&x=2.
\end{cases}
\]
¿Para qué valor de $k$ es continua $f$ en $x=2$?}
\begin{alternativas}
  \item $k=0$.
  \item $k=4$.
  \item $k=2$.
  \item No existe un valor de $k$ que haga continua a $f$ en $x=2$.
\end{alternativas}

\reactivo{6}{Una función $f$ es continua en $[0,2]$ y diferenciable en $(0,2)$.
Se sabe que $f(0)=-1$, $f(2)=3$ y $f'(x)>0$ para todo $x\in(0,2)$.
¿Qué se puede concluir sobre las soluciones de $f(x)=0$ en $(0,2)$?}
\begin{alternativas}
\item No existe ninguna solución.
    \item Existe al menos una solución, pero las condiciones
    no permiten determinar si es única.
    \item Existen al menos dos soluciones distintas.
    \item Existe exactamente una solución.
\end{alternativas}

\reactivo{7}{Sea $f$ diferenciable en un intervalo abierto $I$. Si $f'(x)\geq0$
para todo $x\in I$, ¿cuál de las siguientes conclusiones está garantizada?}
\begin{alternativas}
\item Para cualesquiera $x_1,x_2\in I$ con $x_1<x_2$,
    se cumple $f(x_1)<f(x_2)$.
    \item $f$ es constante en $I$.
    \item $f$ no puede ser constante en ningún subintervalo de $I$.
    \item Para cualesquiera $x_1,x_2\in I$ con $x_1<x_2$,
    se cumple $f(x_1)\leq f(x_2)$.
\end{alternativas}

\reactivo{8}{La ecuación $x^2+xy+y^2=3$ define localmente a $y$ como una
función diferenciable de $x$ alrededor de $(1,1)$. ¿Cuál es la pendiente de
la recta tangente en ese punto?}
\begin{alternativas}
  \item $1$.
  \item $-\dfrac{1}{2}$.
  \item $-1$.
  \item $0$.
\end{alternativas}

\reactivo{9}{Una función $f$ es continua en $[0,3]$ y diferenciable en $(0,3)$.
Se sabe que $f'(x)<0$ para $0<x<1$ y que $f'(x)>0$ para $1<x<3$.
Sin conocer una fórmula para $f$, ¿cuál de las siguientes afirmaciones
es necesariamente verdadera?}
\begin{alternativas}
\item $f$ alcanza su máximo absoluto en $x=1$.
    \item $f$ alcanza su mínimo absoluto en $x=1$.
    \item $f(0)<f(1)$.
    \item $f(3)<f(1)$.
\end{alternativas}

\reactivo{10}{¿Cuál es el polinomio de Taylor de grado dos de $f(x)=e^x$
centrado en $x=0$?}
\begin{alternativas}
  \item $1+x+x^2$.
  \item $x+\dfrac{x^2}{2}$.
  \item $1+\dfrac{x^2}{2}$.
  \item $1+x+\dfrac{x^2}{2}$.
\end{alternativas}

\ifsoluciones
\clearpage
\section*{Clave y correspondencia con el programa}

\begin{center}
\begin{tabular}{ccl}
\hline
Reactivo & Respuesta & Contenido del programa (página 2)\\
\hline
1 & B & 1.1: Supremo e ínfimo de conjuntos acotados\\
2 & D & 2.1--2.2: Convergencia y álgebra de límites\\
3 & A & 2.5: Criterio de Cauchy\\
4 & C & 6.2 y 6.6: Exponencial y regla de L'Hópital\\
5 & B & 3.1 y 3.3: Límites y continuidad\\
6 & D & 4.3 y 5.6: Existencia y unicidad de un cero\\
7 & D & 5.6: Valor medio y monotonía\\
8 & C & 5.1 y 5.5: Tangente y derivación implícita\\
9 & B & 6.3: Extremos de una función\\
10 & D & 6.7: Fórmula de Taylor\\
\hline
\end{tabular}
\end{center}

\section*{Justificaciones}
\begin{enumerate}
\item \textbf{B.} Para $n=1$ se obtiene $0$, el mínimo. Todos los elementos
son menores que $1$ y convergen a $1$, de modo que $\sup S=1$.
Como $1\notin S$, no existe máximo.

\item \textbf{D.} Dividiendo numerador y denominador por $n^2$:
\[
a_n=\frac{3+1/n^2}{2+1/n}\longrightarrow\frac{3}{2}.
\]

\item \textbf{A.} La definición exige que cualesquiera dos términos de una
cola de la sucesión estén arbitrariamente próximos. B expresa acotamiento
de las diferencias, C expresa convergencia a cero y D exige una sucesión
eventualmente constante; ninguna de estas tres condiciones es la definición
de Cauchy.

\item \textbf{C.} Se presenta la forma indeterminada $0/0$. La regla de
L'Hópital es aplicable: las funciones son diferenciables en un entorno de
cero y la derivada del denominador es $1\neq0$. Por tanto,
\[
\lim_{x\to0}\frac{e^{2x}-1}{x}
=\lim_{x\to0}2e^{2x}=2.
\]

\item \textbf{B.} Para $x\neq2$, $f(x)=x+2$. Así,
$\lim_{x\to2}f(x)=4$; la continuidad exige que $f(2)=k=4$.

\item \textbf{D.} La continuidad y el cambio de signo
garantizan al menos un cero interior. Como $f'>0$, el teorema del valor
medio implica que $f$ es estrictamente creciente y no puede tomar el
valor $0$ en dos puntos distintos. Así, el cero existe y es único.

\item \textbf{D.} Para $x_1<x_2$, el teorema del valor medio
da $f(x_2)-f(x_1)=f'(c)(x_2-x_1)\geq0$ para algún $c\in(x_1,x_2)$.
Esto garantiza que $f$ es no decreciente. No garantiza crecimiento
estricto: una función constante satisface la hipótesis. Tampoco obliga
a que $f$ sea constante, como muestra $f(x)=x$.

\item \textbf{C.} Derivando implícitamente:
\[
2x+y+xy'+2yy'=0,
\qquad y'=-\frac{2x+y}{x+2y}.
\]
En $(1,1)$, el denominador vale $3$ y la pendiente es $-3/3=-1$.

\item \textbf{B.} El signo negativo de $f'$ implica que $f$
es estrictamente decreciente en $[0,1]$; el signo positivo implica que es
estrictamente creciente en $[1,3]$. Ambas afirmaciones se justifican
aplicando el teorema del valor medio a pares de puntos de cada intervalo.
Por tanto, $f(1)<f(x)$ para todo $x\in[0,3]$ distinto de $1$.

\item \textbf{D.} Puesto que $f(0)=f'(0)=f''(0)=1$,
\[
P_2(x)=f(0)+f'(0)x+\frac{f''(0)}{2!}x^2
=1+x+\frac{x^2}{2}.
\]
\end{enumerate}
\fi
\end{document}
